\documentclass[11pt,a4paper]{article}
\usepackage[T1]{fontenc}
\usepackage{lmodern,amsmath,amssymb,amsthm,mathtools,microtype}
\usepackage[margin=27mm]{geometry}
\usepackage[colorlinks=true,linkcolor=blue,citecolor=blue,urlcolor=blue]{hyperref}
\newtheorem{theorem}{Theorem}[section]
\newtheorem{lemma}[theorem]{Lemma}
\newtheorem{corollary}[theorem]{Corollary}
\newcommand{\Q}{\mathbb Q}\newcommand{\Z}{\mathbb Z}
\setlength{\emergencystretch}{2em}
\title{Effective heights of elliptic sums with torsion translates}
\author{}\date{17 November 2026}
\begin{document}\maketitle

\begin{abstract}
For a rank-one family on $y^2=x^3-d^2x$ over a totally real
number field, an integral translation of an abscissa sum has height
at least the height of any summand, up to an explicit logarithmic
index loss and a ramification constant. Arbitrary torsion translates,
negative indices, repeated points and finite points of index zero
are allowed. Over $\Q(\sqrt2)$ the ramification constant vanishes.
On the normalized elliptic source this gives an effective quadratic
index bound on the entire finite point domain.
\end{abstract}

\section{The uniform height inequality}

Let $K$ be totally real, let $k=[K:\Q]$, and let
$d\in\mathcal O_K$ be totally positive. Put
\[
 E:\quad y^2=x^3-d^2x.
\]
Fix an infinite-order point $P\in E(K)$ and an integer $r\ge1$.
For $1\le i\le r$ let
\[
 Q_i=[n_i]P+T_i\in E(K)\setminus\{O\},\qquad
 n_i\in\Z,\qquad T_i\in E(K)_{\mathrm{tors}}.
\]
Every $Q_i$ is finite, but its index may be zero. Set
\[
 M=\max(1,|n_1|,\ldots,|n_r|),\qquad
 Z=\sum_{i=1}^r x(Q_i)+rd+1.
\]
We use the absolute logarithmic Weil height, with finite valuation
$v_\mathfrak p$ normalized to have value group $\Z$ and weight
$f_\mathfrak p\log p/k$.
Define
\[
 A_r=(r-1)2^{r-1},\qquad
 B_r=\sum_{j=1}^r\binom{r-1}{j-1}\log j,
\]
and the explicit ramification allowance
\[
 R_K=\frac2k\sum_{\substack{\mathfrak p\mid p\\p\ne2}}
       f_\mathfrak p\left\lfloor\frac{e_\mathfrak p-1}{2}\right\rfloor\log p.
\]
Only odd primes with ramification index at least three contribute.

\begin{theorem}\label{thm:main}
With the preceding notation,
\[
 h(Z)\ge\max_i h(x(Q_i))-A_r\log M-B_r-h(d)-R_K.
\]
Moreover
\[
 R_K\le\frac1k\log|\operatorname{disc}(K)|.
\]
In particular $R_K=0$ when every odd ramification index is at most
two, including every totally real multiquadratic field.
\end{theorem}

The result bounds a single algebraic abscissa sum. It makes no
claim that its field trace has the same effective lower bound.
Its constants do not depend on the chosen torsion translates or
on an upper bound for their orders.

\section{Local principal parts and ramification}

Write $D=d^2$ and $z=-x/y$ at the origin. The formal
expansions are
\begin{align}
 x(z)&=z^{-2}+Dz^2-D^2z^6+2D^3z^{10}-\cdots,\label{eq:x}\\
 \ell(z)&=\sum_{j\ge0}
 \frac{(-1)^j\binom{2j}{j}D^j}{4j+1}z^{4j+1}.
 \label{eq:log}
\end{align}
The first expansion follows by solving the quadratic equation for
$x$ in terms of $z$. The second follows by integrating the invariant
differential
$dx/(2y)=(1+4Dz^4)^{-1/2}dz$.
These computations and the positive-subset argument below extend
the integral principal-part method of~\cite{Sums}.

Fix a finite place $\mathfrak p\mid p$ with ramification index $e$.
A nonidentity formal point has
\[
 a=v_\mathfrak p(z)>0,\qquad v_\mathfrak p(x)=-2a.
\]
Call it deep if $a\ge\lceil e/2\rceil$ for odd $p$, and if
$a\ge1$ for $p=2$.
Let
\[
 b_\mathfrak p=\begin{cases}2\lfloor(e-1)/2\rfloor,&p\ne2,\\0,&p=2.\end{cases}
\]
A finite point which is not deep has
$v_\mathfrak p(x)\ge-b_\mathfrak p$.

\begin{lemma}\label{lem:principal}
On the deep formal subgroup,
\[
 v_\mathfrak p(\ell(z))=v_\mathfrak p(z),\qquad
 x(z)-\ell(z)^{-2}\in\mathcal O_{K_\mathfrak p}.
\]
The logarithm is injective there.
\end{lemma}
\begin{proof}
For odd $p$ and $j\ge1$,
$v_p(4j+1)\le2j-1$, since $3^{2j}>4j+1$.
If $2a\ge e$, the $j$th nonconstant
term of~\eqref{eq:log} has valuation at least
\[
 (4j+1)a-e\,v_p(4j+1)\ge(4j+1)a-2a(2j-1)=3a.
\]
The valuations tend to infinity, since $v_p(4j+1)=O(\log j)$.
For $p=2$ every denominator $4j+1$ is a unit and the same terms
have valuation at least $5a$. Hence in both cases
$\ell(z)=z(1+\eta)$ with $v_\mathfrak p(\eta)\ge2a$.
It follows that $\ell(z)^{-2}-z^{-2}$ is integral.
The terms of~\eqref{eq:x} after $z^{-2}$ are integral as well.
The formal logarithm is a group homomorphism, and the valuation
identity proves that its kernel on this subgroup is trivial.
\end{proof}

The deep subgroup is open in the compact group $E(K_\mathfrak p)$.
Choose $m>0$ with $[m]P$ in it and set $L=\ell([m]P)/m$.
The homomorphism property makes this independent of $m$.
It is nonzero: otherwise injectivity would make $P$ torsion.
If $Q_i$ is deep, multiplication by a common multiple of $m$ and
the order of $T_i$, followed by the homomorphism identity, gives
\[
 \ell(Q_i)=n_iL.
\]
In particular $n_i\ne0$, and
\[
 x(Q_i)=(n_iL)^{-2}+\varepsilon_i,
 \qquad v_\mathfrak p(\varepsilon_i)\ge0.
\]
No assertion of integrality for a shallow torsion point is needed.

\section{The positive subset numerators}

Fix $i$, and suppose first that $Q_i$ is deep at $\mathfrak p$.
Let $J$ be the set of all deep indices at that place and put
\[
 C_J=\sum_{j\in J}\prod_{a\in J\setminus\{j\}}n_a^2
       \in\Z_{>0}.
\]
Adding the preceding expansions and the other abscissas gives
\[
 Z=(n_iL)^{-2}
      \frac{C_J}{\prod_{j\in J\setminus\{i\}}n_j^2}
      +\varepsilon,
 \qquad v_\mathfrak p(\varepsilon)\ge-b_\mathfrak p.
\]
The integral translation $rd+1$ is included in $\varepsilon$.
Let $H=-v_\mathfrak p(x(Q_i))$ and let $v(A)$ denote the
valuation of the displayed rational quotient.
If $H-v(A)>b_\mathfrak p$, its principal part dominates and
the loss of local height is $v(A)$. Otherwise the loss is at
most $H\le v(A)+b_\mathfrak p$. Since the denominator of $A$
is a product of integer squares, $v(A)\le v(C_J)$.
Thus in all cases the local loss is at most
\[
 v_\mathfrak p(C_J)+b_\mathfrak p.
\]
If $Q_i$ is not deep, its entire local height is already at most
$b_\mathfrak p$, giving the same bound without a subset term.

For fixed $i$, summing over finite places is bounded by summing
$\log C_J$ over all subsets containing $i$ and consisting of
nonzero indices, plus $R_K$. For $|J|=j$,
\[
 C_J\le jM^{2(j-1)}.
\]
There are at most $\binom{r-1}{j-1}$ such subsets.
Consequently their total logarithmic cost is at most
$A_r\log M+B_r$. If $n_i=0$, its point is never deep and the
same upper bound remains valid.

At a real embedding, $x(x-d)(x+d)\ge0$ implies $x\ge-d$.
Hence $Z\ge1$. A nonnegative summand is at most $Z$, and a
negative summand has absolute value at most $d$. Therefore
\[
 \max(1,|x(Q_i)|)\le\max(1,d)\max(1,Z).
\]
The sum of these archimedean losses is $h(d)$ because $d$ is
integral. Combining the finite and infinite estimates proves the
first assertion of Theorem~\ref{thm:main}.
Finally $2\lfloor(e-1)/2\rfloor\le e-1$, and the exponent of
the local different is at least $e-1$. The discriminant formula
bounds $R_K$ by $\log|\operatorname{disc}(K)|/k$.
At an odd prime the inertia subgroup in a multiquadratic extension
is tame and cyclic of exponent at most two; hence its ramification
index is at most two, proving the last assertion.

\section{The complete normalized point domain}

Take $K=\Q(\sqrt2)$, $d=3-\sqrt2$ and
\[
 R_0=(1+\sqrt2,2\sqrt2),\qquad P_0=2R_0.
\]
The retained normalized elliptic source~\cite{Source} supplies
\[
 E(K)=\Z R_0\oplus E[2](K),\qquad
 \widehat h(P_0)\ge\tfrac12,
\]
and its uniform duplication inequality
\[
 h(x(2Q))\le4h(x(Q))+\log676
\]
for non-torsion $Q$. We identify these as arithmetic inputs to
the following numerical corollary; Theorem~\ref{thm:main} itself
requires none of this particular group computation.

\begin{corollary}\label{cor:four}
For four finite points $Q_i=[m_i]R_0+T_i$, with $T_i\in E[2](K)$,
put $Z=\sum_i x(Q_i)+4d+1$ and
\[
 B=\max(1,|m_1|/2,|m_2|,|m_3|,|2m_4+1|).
\]
Then
\[
                         B^2<40h(Z)+2285.
\]
\end{corollary}
\begin{proof}
The odd primes of $K$ are unramified, $h(d)=\tfrac12\log7$,
$A_4=24$ and $B_4=\log864$.
Iterating the duplication inequality and using the canonical-height
limit gives
\[
 h(x(Q_i))\ge2\widehat h(Q_i)-\tfrac13\log676
             \ge m_i^2/4-3
\]
when $m_i\ne0$. Let $M=\max(1,|m_i|)$.
If every index is zero the asserted bounds below are immediate.
Otherwise apply Theorem~\ref{thm:main} to an index attaining $M$,
use the weaker constant six in the last inequality, and use
\[
 24\log M-M^2/8\le12\log96-12<43.
\]
Since $6+43+\log864+\tfrac12\log7<57$, this yields
$M^2<8h(Z)+456$.
Finally $B\le2M+1$ and $(2M+1)^2\le5M^2+5$.
\end{proof}

\begin{thebibliography}{9}
\bibitem{Sums}
\emph{Effective heights of sums of elliptic abscissas},
8 November 2026,
\nolinkurl{papers/effective_sums/effective_elliptic_sums_20261108_v10.tex}.
\bibitem{Source}
\(\Z\Q\) Diophantine Project contributors,
\emph{$\Z$ is Diophantine in $\Q$}, Version 19,
17 November 2026; the normalized duplication bounds preceding
equation \textup{(o8-indexheight)}, and the half-generator domain.
Pinned commit \texttt{498efba6fe75e9d28d6804f1391d7f70e4150afa},
\url{https://github.com/elliotglazer/diophantine-z-q}.
\end{thebibliography}
\end{document}
